\documentclass[11pt,a4paper]{article}
\usepackage[margin=24mm]{geometry}
\usepackage{amsmath,amssymb,booktabs,longtable,array,hyperref}
\hypersetup{hidelinks}
\emergencystretch=2em
\newcommand{\Lc}{\textbf{(c)}}
\newcommand{\Ld}{\textbf{(d)}}
\newcommand{\La}{\textbf{(a)}}
\newcommand{\G}{\texttt{det\_geometry.py}}
\newcommand{\Scheck}{\texttt{det\_spectrum.py}}
\newcommand{\Ccheck}{\texttt{det\_charges.py}}
\newcommand{\dd}{\mathrm d}
\newcommand{\R}{\mathbb R}
\title{The projected one-loop Boussinesq determinant:\\
potential endpoints, collective measure, and a conditional dispersion invariant}
\author{Research note for Gustav Delius}
\date{4 November 2026}
\begin{document}
\maketitle
\begin{abstract}
\Lc\ We construct the fourth-order potential fluctuation operator, determine
its bounded static zero modes in a specified decaying-derivative class, and
calculate the collective symplectic form including the shift charge. Translation,
velocity and shift alone give an odd-dimensional, rank-two form. An explicit
auxiliary momentum direction completes it and gives a checked finite measure
and projector. This is an instantaneous phase-space construction, not a
functional determinant. The shift and width tangents have divergent
configuration-space kinetic norms; the enlarged phase-space construction does
not remove that temporal-domain problem.

\Lc\ Conditional on the inherited continuum trace prescription, both published
charge pairs give the same correction at fixed momentum:
$E_1|_P=\Delta E_{\rm mode}-v\Delta P_{\rm mode}=-L_1$.
Its small-amplitude coefficient is $E_1|_P/|P_{\rm cl}|\to-2/\pi$.
\Ld\ Neither this note nor the inherited trace establishes that this is the
absolutely normalized projected determinant. The obstruction is specified as
a missing common boundary/constraint evolution problem, followed by spectral
completeness and the phase of its time-boundary determinant. Thus the outcome
is the checked partial-result alternative in the task, not a claim of a
completed one-soliton dispersion or two-soliton S-matrix.
\end{abstract}

\section{Scope, evidence and conventions}
Labels mean literature theorem (a), literature heuristic (b), project
derivation (c), and guess or unestablished identification (d).
Each section names its checking script; all scripts are new, independent of
imports from earlier project calculations, and live in \texttt{phase2/scripts/}.
Every displayed derived formula below is covered by those scripts, using
exact SymPy identities or 55--65-digit mpmath tests. Numerical residual tests
are evidence at the listed points, not proofs of completeness.

The input is the good equation and the specified normal ordering. No new
regularization or physical endpoint prescription is selected. The symmetric
cutoff below is only a diagnostic of the already published conditional
equal-mode-count expression. The auxiliary finite-dimensional completion is
also a diagnostic, not a selection of quantum boundary degrees of freedom.
No Phase~2i formula or data is used. No existing file is changed.

\Lc\ (\G) It is useful to keep the shift coordinate explicitly:
\begin{align}
 \xi&=x-X-vt,\qquad k=\sqrt{1-v^2}=\sin\alpha,\quad v=\epsilon\cos\alpha,
 \quad 0<\alpha<\pi/2,\\
 \Phi(\xi;v,C)&=C+k\tanh(k\xi/2),\qquad
 u_0=\Phi'=\frac{k^2}{2}\operatorname{sech}^2(k\xi/2),\quad \pi_0=-vu_0,
 \label{eq:profile}\\
 M&=2k,\qquad \Pi=\int\pi_0\dd x=-2vk,\qquad
 P=-\int\pi_0u_0\dd x=\frac23vk^3,\\
 E_0&=\frac23k^3-\frac8{15}k^5,\qquad
 L_0=vP-E_0=-\frac2{15}k^5,\qquad \frac{\dd L_0}{\dd v}=P.
\end{align}
This is a coordinate dictionary, not a change of normalization: the old
left-endpoint coordinate is $C_L=C-k$, so $\Phi=C_L+k(1+\tanh(k\xi/2))$.
Keeping $C$ free avoids silently identifying a physical shift symmetry with
a gauge choice. The fixed symplectic convention remains
$\Omega=\int\dd\pi\wedge\dd\phi$.

\section{Operator and spatial boundary problem}
\paragraph{Quadratic operator. \Lc\ (\G, \Scheck).}
For $\phi=\Phi+\eta$ the lab-frame field equation and quadratic action give
\begin{align}
 \phi_{tt}-\phi_{xx}+\phi_{xxxx}+3(\phi_x^2)_x&=0,\\
 S_2&=\frac12\int\dd t\dd\xi\,
 \{(\eta_t-v\eta_\xi)^2-(1-6u_0)\eta_\xi^2-\eta_{\xi\xi}^2\},\\
 \mathcal D_v&=(\partial_t-v\partial_\xi)^2+\mathcal L,
 \qquad \mathcal L=\partial_\xi^4-\partial_\xi^2
                  +6\partial_\xi(u_0\partial_\xi),\\
 \mathcal K_v&=\mathcal L+v^2\partial_\xi^2
        =-\partial_\xi\mathcal H_v\partial_\xi,
 \qquad \mathcal H_v=-\partial_\xi^2+k^2-6u_0,\label{eq:K}\\
 \mathcal D_v(e^{-i\nu t}f)&=e^{-i\nu t}
       (\mathcal K_v+2iv\nu\partial_\xi-\nu^2)f.
\end{align}
The Hessian in $S_2=-\frac12\int\eta\mathcal D_v\eta$ has the corresponding
overall minus sign after integration by parts. In particular the mixed
co-moving derivative has the plus sign displayed above.
The potential operator is not the $u$-fluctuation operator.

\paragraph{Boundary form and admissible variations. \Lc\ (\G).}
Writing $A=-u+3u^2+u_{xx}$, the spatial boundary term of the first variation
of the original action is
\begin{equation}
 \delta S\big|_{\rm space}=\int\dd t\,[A\delta\phi-u_x\delta\phi_x]_{-\infty}^{+\infty}.
 \label{eq:variation}
\end{equation}
For the classical pulse these coefficients tend to zero. Thus finite,
time-independent endpoint variations are not excluded by the classical
spatial variational principle alone. The Green identity is
\begin{align}
 f\mathcal K_vg-g\mathcal K_vf&=\partial_\xi B_v(f,g),\\
 B_v(f,g)&=fg'''-f'g''+f''g'-f'''g
       +(6u_0+v^2-1)(fg'-gf').\label{eq:green}
\end{align}
Bounded $f,g$ with decaying derivatives have zero flux. Clamped fluctuations
$f=f'=0$ at both ends of a finite interval also make it vanish; this is a
mathematical example of a domain, not the selected physical regulator.
Neither zero boundary flux nor formal symmetry alone defines a determinant.

\Lc\ (\G) Two different linear spaces must be distinguished. For a fixed
potential-endpoint sector one may use localized smooth initial data
$(\eta,\rho)$ with $\eta(\pm\infty)=0$. For the following finite pairings an
enlarged test space is
\begin{equation}
 \eta=c+m\tanh(k\xi/2)+s(\xi),\qquad
 \rho=r(\xi),\qquad s,r\in\mathcal S(\R).
 \label{eq:testspace}
\end{equation}
All pairings converge on this space; it includes the bounded collective tails.
This is not a claim that it is the Hilbert space of the Gaussian integral.
For smooth finite-energy histories with uniform endpoint limits and
$\pi(\pm\infty)=0$, local conservation gives
\begin{equation}
 \dot M=[\pi]_{-\infty}^{+\infty}=0,\qquad
 \dot\Pi=[u-3u^2-u_{xx}]_{-\infty}^{+\infty}=0.\label{eq:charges}
\end{equation}
The assumptions on limits matter: arbitrary $L^2$ functions need not have
pointwise limits. These statements concern the smooth decaying class here.

\paragraph{The temporal obstruction made quantitative. \Lc\ (\G).}
In the centred coordinates of \eqref{eq:profile},
\begin{equation}
 \Phi_v=-\frac vk\left[\tanh(k\xi/2)
                  +\frac{k\xi}{2}\operatorname{sech}^2(k\xi/2)\right],
 \quad (\Phi_v)_\pm=\mp\frac vk.
\end{equation}
Consequently, in a symmetric box $[-R,R]$,
\begin{align}
 \|1\|^2&=2R,\qquad
 \|\Phi_v\|^2=2R\frac{v^2}{k^2}+O(1),\\
 \frac12\int_{-R}^R\dot\Phi^2\dd\xi
  &=R\left(\dot C^2+\frac{v^2}{k^2}\dot v^2\right)+O(1)
  \quad\hbox{for a time-modulated profile.}\label{eq:norm}
\end{align}
The bounded remainder includes translation and other localized terms.
At fixed nonzero $v$, these independent collective paths do not have a finite
configuration-space kinetic term as $R\to\infty$. A finite-action history
can still have changing coordinates if the complementary field has exactly
compensating endpoint tails. Such a compensation couples the collective and
field constraints; it cannot be omitted when using the instantaneous measure.
This does not forbid comparing stationary solitons in different endpoint
sectors, nor prove that no quantum soliton theory can be defined.

\section{Zero, collective and nonzero modes}
\paragraph{Static kernel and generalized mode. \Lc\ (\Scheck).}
In the bounded, decaying-derivative class just described,
\begin{equation}
 \mathcal K_v1=0,\qquad \mathcal K_vu_0=0,\qquad
 \mathcal K_v\Phi_v=-2vu_0',\qquad
 \eta_v(\xi,t)=\Phi_v(\xi)-t u_0(\xi),\quad
 \mathcal D_v\eta_v=0.\label{eq:jacobi}
\end{equation}
Translation is $-u_0$, shift is $1$, and the last field is the velocity
Jacobi solution at fixed lab coordinates. Width variation is not a symmetry
zero mode of $\mathcal K_v$ for $v\ne0$. Its time-secular term is essential.
The formal shift partner $\eta=t$ also solves the equation, but has uniform
$\delta\pi=1$ and hence infinite kinetic energy on the line; it is excluded
from the finite-energy class. It cannot supply the missing finite shift-charge
direction in the collective measure.

\Lc\ (\G, \Scheck) This static kernel is exhaustive in the stated class:
integrating $\mathcal K_vf=0$ gives $\mathcal H_vf'=0$, since derivatives
decay and hence the integration constant is zero. Direct substitution gives
$\mathcal H_vu_0'=0$. Any other solution with decaying derivative has constant
Schr\"odinger Wronskian with $u_0'$, which is zero at infinity, and is therefore
proportional to it by uniqueness of the second-order ODE. Integration gives
$f=c_0+c_1u_0$. In the $L^2$ potential domain only translation remains;
in a fixed-endpoint domain the shift and $\eta_v$ are excluded.
This argument does not enumerate nonzero-frequency internal modes or establish
completeness of the frequency-dependent pencil.

\paragraph{Continuum checks and a heavy-branch correction. \Lc\ (\Scheck).}
For real $q\ne0$, $\omega=|q|\sqrt{1+q^2}$ and $\nu=\omega-vq$.
The vacuum pencil is
\begin{equation}
 q^4+k^2q^2-2v\nu q-\nu^2=0,\qquad
 \partial_q^2(\hbox{left side})=12q^2+2k^2>0.
\end{equation}
It has two real roots, one of each sign, and a non-real conjugate pair for
real nonzero $\nu$. A convenient exact continuum solution is
\begin{align}
 f_q(\xi)&=\partial_\xi\left[
 e^{iq\xi}\frac{1+A e^{k\xi}}{1+e^{k\xi}}\right],\\
 A&=\frac{(p_a-q_a)(p_b-q_b)}{(p_a-q_b)(p_b-q_a)},\qquad
 p_{a,b}=\frac{v/\sqrt3\pm k}{2},\quad
 q_{a,b}=\frac{\operatorname{sgn}(q)\sqrt{1+q^2}}{2\sqrt3}\pm\frac{iq}{2}.
\end{align}
The checked pencil residual is below $9\times10^{-61}$ over the tested
Jacobi, continuum and discrete modes. The same ratio gives $|A|=1$;
the asymptotic amplitudes of $f_q$ are $iq$ and $iqA$.

\Lc\ (\Scheck) On the heavy branch $\pi/3<\alpha<\pi/2$ define
\begin{align}
 K_-&=\sin(\alpha-\pi/3),\quad K_+=\sin(\alpha+\pi/3),\quad K_-+K_+=k,\\
 \gamma&=\sqrt3(k^2-3/4),\qquad
 f_\pm^{\rm sp}(\xi)=\partial_\xi\frac{e^{K_\pm\xi}}{1+e^{k\xi}}.
\end{align}
The two localized solutions are
\begin{equation}
 e^{\epsilon\gamma t}f_-^{\rm sp}(\xi),\qquad
 e^{-\epsilon\gamma t}f_+^{\rm sp}(\xi).\label{eq:heavy}
\end{equation}
The same spatial profile cannot be used for both rates when $v\ne0$.
The shorthand in \texttt{oneloop\_phase.tex} does that; its opposite-rate
potential-field residual at $\alpha=1.2$, $\xi=0.37$ is
$0.00733252608902685$ in this normalization, for either orientation.
Using the complementary $K$ fixes the residual without a tolerance change.
For the light branch the displayed $K_-$ mode is not localized.
\Ld\ A real, stable-particle dispersion on the heavy branch is not supplied
by an oscillatory Gaussian trace in the presence of these growing modes.
The conditional real-energy benchmark below is restricted to
$0<\alpha<\pi/3$; this is a limitation of that benchmark, not an unannounced
change of equation. The cusp requires separate treatment.

\section{Collective symplectic split and what its Jacobian means}
\paragraph{Three tangents are insufficient. \Lc\ (\G).}
At $t=0$ order the exact soliton-family tangents as $(v,X,C)$. Integration gives
\begin{equation}
 W_3=\begin{pmatrix}0&P_v&\Pi_v\\-P_v&0&0\\-\Pi_v&0&0\end{pmatrix},\qquad
 P_v=\frac23k(1-4v^2),\quad \Pi_v=\frac{2(2v^2-1)}k.
 \label{eq:W3}
\end{equation}
It has rank two and null vector
$\Pi_v\partial_X-P_v\partial_C$. This degeneracy of the restricted
solution family is not evidence that the shift is a gauge redundancy of the
full field theory. The shift pairs with its charge $\Pi$, which is already
tied to $v$ on that family. Inverting a three-mode Wronskian is impossible.

\paragraph{An explicit completion and finite measure. \Lc\ (\G).}
To expose the missing direction without choosing a physical boundary theory,
introduce at one time an auxiliary variable $r$ by
\begin{equation}
 \pi=-vu_0+r w,\qquad
 w=\frac3k u_0-\frac{15}{2k^3}u_0^2,\qquad
 \int w\dd\xi=1,\quad \int wu_0\dd\xi=0.
\end{equation}
These follow from $\int u_0^n\dd\xi=2k,2k^3/3,4k^5/15$ for $n=1,2,3$.
The vector $Z_r=(0,w)$ is localized, changes $\Pi$ independently, and does
not change $P$ to first order. It is not a stationary soliton tangent and is
not to be removed as another solution zero mode. At $r=0$, parity gives
$\int w\Phi_v=0$. In the order $(v,X,C,r)$ the completed matrix and density are
\begin{equation}
 W=\begin{pmatrix}
 0&P_v&\Pi_v&0\\-P_v&0&0&0\\-\Pi_v&0&0&-1\\0&0&1&0
 \end{pmatrix},\qquad
 \det W=P_v^2,\qquad
 \dd\mu_{\rm coll}=|P_v|\dd v\dd X\dd C\dd r.
 \label{eq:W4}
\end{equation}
This density is local at $r=0$ and omits universal phase-space normalization
constants; it is not an absolute state normalization. The quadrature checks
at $v=0.8,-0.7,0.3$ have maximal error below $2.1\times10^{-56}$.
The completion fails at $P_v=0$, namely $|v|=1/2$; it does not resolve the cusp.

\paragraph{Projector, constraints and Jacobian. \Lc\ (\G).}
For the completed tangent map $Z$ and $P_v\ne0$, define
\begin{align}
 b_A&=\Omega(Z_A,\eta),\qquad
 \mathcal P\eta=Z_A(W^{-1})^{AB}b_B,\qquad \mathcal Q=1-\mathcal P,\\
 \mathcal P^2&=\mathcal P,\qquad \Omega(Z_A,\mathcal Q\eta)=0,\\
 \Omega\big|_{\rm tangent}&=\tfrac12W_{AB}\dd z^A\wedge\dd z^B+\Omega_\perp.
\end{align}
The last line denotes the block decomposition of the two-form.
For any finite symplectic realization, the associated
change of Liouville measure has Jacobian $\sqrt{\det W}=|P_v|$ times the
transverse measure. The exact finite constraint identity is
\begin{equation}
 1=|\det W|\int\dd^4c\,
       \delta^{(4)}\!\left(\Omega(Z,\eta-Zc)\right).\label{eq:FP}
\end{equation}
Thus the constraint determinant is $|\det W|$, whereas the collective
symplectic density is $\sqrt{\det W}$; these must not be multiplied as if
they were two independent Jacobians. The script checks a six-dimensional
canonical embedding, the projector, the cross block, and the full change of
variables determinant exactly. Equations \eqref{eq:W4}--\eqref{eq:FP} apply
to the tangent Gaussian geometry, not to an unconstructed infinite product.

\paragraph{The remaining determinant problem. \Lc/\Ld.}
\Lc\ (\G) The projector preserves the enlarged test space
\eqref{eq:testspace}; in general it does not preserve zero endpoint
variations. It can create both constant and opposite-sign tails. Moreover
$Z_r$ is an auxiliary initial-data direction, so using it through time requires
an evolved constraint basis, or explicit $\dot{\mathcal Q}$ terms. An
instantaneous block diagonalization of $\Omega$ is not a block
diagonalization of the quadratic action. Already for a finite symmetric block
matrix the determinant is
\begin{equation}
 \det\begin{pmatrix}A&B\\B^T&C\end{pmatrix}
 =\det C\,\det(A-BC^{-1}B^T),\qquad C\text{ invertible}.\label{eq:schur}
\end{equation}
The script checks the scalar nontrivial case; the general identity follows
by triangular elimination. The Schur term is lost if the mixed constraints
or the time-dependent basis are silently dropped.

\Lc\ (\G) Static or Jacobi solutions also need not be finite-time determinant
zero modes. With temporal Dirichlet endpoints, the elementary second
difference matrix has determinant $N+1$, despite its continuum equation
having solutions $1,t$. An analogous decision about temporal endpoints is
necessary before placing a prime on the Boussinesq determinant.

\Ld\ The precise failed step is therefore promotion of
$|P_v|\dd v\dd X\dd C\dd r$ to a factorized Gaussian path measure in which
$v(t),C(t)$ and a fixed-endpoint radiation field vary independently. Equation
\eqref{eq:norm} excludes that promotion on the infinite-line kinetic domain.
Allowing compensating radiation tails, reducing at specified endpoint/shift
charges, or adding boundary variables are different possible constructions;
none is selected here. One must specify which construction defines the
quantum amplitude, propagate its constraints, then evaluate the resulting
Schur complement and time-endpoint determinant. This note supplies neither
that determinant nor the nonzero-frequency completeness theorem. No inference
of zero relative Maslov phase is made from the finite measure.

\section{Specified normal ordering and the conditional trace}
\paragraph{Counterterm. \Lc\ (\Ccheck).}
With vacuum frequency $\omega(q)=|q|\sqrt{1+q^2}$, the symmetric-cutoff vacuum
contraction and the prescribed Wick polynomial are
\begin{equation}
 C_\Lambda=\langle\phi_x^2\rangle_{0,\Lambda}
 =\frac{\hbar}{2}\int_{-\Lambda}^{\Lambda}\frac{\dd q}{2\pi}
                       \frac{q^2}{\omega(q)}
 =\frac{\hbar}{2\pi}(\sqrt{1+\Lambda^2}-1),\quad
 :\!\phi_x^3\!:\,=\phi_x^3-3C_\Lambda\phi_x.
\end{equation}
Consequently $\delta L_{\rm ct}=-3C_\Lambda M$ and
$\delta E_{\rm ct}=+3C_\Lambda M$ on a fixed background.
The script verifies the Gaussian moment identity and the elementary integral.
This is the specified vacuum normal ordering, with no additional finite
multiple of $M$ chosen. It leaves the bulk classical equation unchanged but
contributes $-3C_\Lambda\delta M$ to the spatial endpoint variation.
\Ld\ Whether that term cancels the corresponding determinant endpoint
variation is a question for the renormalized boundary problem; the bare
counterterm alone is neither proof of an anomaly nor a Ward identity.

\paragraph{Conditional benchmark. \Lc\ (\Ccheck).}
For $0<\alpha<\pi/3$ put $\varphi_r=(\alpha+2\pi/3,\alpha-2\pi/3,\alpha)$,
$c_r=(1,1,-2)$. The inherited continuum phase has derivative
\begin{equation}
 \theta'(q)=\sum_r\frac{c_r\sin\varphi_r}{q^2+\sin^2\varphi_r}
 \left(1+\frac{\epsilon\operatorname{sgn}(q)\cos\varphi_r}{\sqrt{1+q^2}}\right).
\end{equation}
On the branches $\theta(\pm\infty)=0$, let $e,p$ denote coefficients of
$\hbar$, not full charge shifts. If the equal-mode-count continuum trace is
the relevant trace, it gives
\begin{align}
 B_\Lambda&=\frac{\Lambda\sqrt{1+\Lambda^2}}{4\pi}
                      [\theta(\Lambda)-\theta(-\Lambda)],\\
 e&=\lim_{\Lambda\to\infty}\left\{
 \frac1{4\pi}\int_{-\Lambda}^{\Lambda}\omega\theta'\dd q-B_\Lambda
             +\frac{3k}{\pi}(\sqrt{1+\Lambda^2}-1)\right\},\\
 p&=\frac1{4\pi}\int_\R q\theta'\dd q,\qquad
 \frac{B_\Lambda}{\Lambda}\longrightarrow\frac{3k}{2\pi},\\
 e(\alpha)&=\frac1{4\pi}\sum_r c_r
       [2\sin\varphi_r-\sin(2\varphi_r)\log|\tan(\varphi_r/2)|],\\
 p(\alpha,\epsilon)&=-\frac\epsilon{2\pi}\sum_r c_r\sin\varphi_r
                                         \log|\tan(\varphi_r/2)|,\label{eq:mode}\\
 L_1(v)&=vp-e.
\end{align}
The derivative of the implemented phase is checked independently, its trace
antiderivative is checked by finite-cutoff quadrature, and the renormalized
limit at $\Lambda=10^{22}$ agrees to below $1.8\times10^{-23}$. This residual
includes the finite-cutoff error and is not claimed to be a 65-digit continuum
limit. Omitting $B_\Lambda$ leaves a linear divergence after the same
normal-ordering subtraction; it is not a harmless finite normalization.

\paragraph{Verdict on Phase 2h. \Lc/\Ld\ (\Ccheck).}
The conditional derivative is indeed
\begin{equation}
 p_{\rm can}=\frac{\dd L_1}{\dd v}
 =-\frac{\epsilon\sqrt3\cos\alpha}{\pi}
 \log\left|\frac{\tan((\alpha+2\pi/3)/2)}{\tan((\alpha-2\pi/3)/2)}\right|,
 \qquad e_{\rm can}=vp_{\rm can}-L_1.\label{eq:canonical}
\end{equation}
This is a checked conditional formula \Lc. Its identification with a
renormalized physical momentum, and identification of $L_1$ with the projected
determinant, remain \Ld. Differentiating across $M(v)$ changes endpoint sector;
an on-shell charge argument requires the quantum boundary variation to be
controlled. The finite Jacobian calculation does not supply that argument.

\section{What can nevertheless be said at fixed physical momentum}
\paragraph{An invariant conditional answer. \Lc\ (\Ccheck).}
For any smooth classical branch with $P_0'\ne0$ and $E_0'=vP_0'$, eliminating
the parameter in $E=E_0+\hbar e$, $P=P_0+\hbar p$ gives
\begin{equation}
 E(P)=E_0(P)+\hbar[e-vp]_{\alpha=\alpha_0(P)}+O(\hbar^2).
\end{equation}
A tangent relabeling changes $(e,p)$ by $(v\,\delta p,\delta p)$ and leaves
this coefficient invariant. Applying it to \eqref{eq:mode} and
\eqref{eq:canonical} gives
\begin{equation}
 \boxed{\quad E_1\big|_P=e-vp=e_{\rm can}-vp_{\rm can}=-L_1.\quad}
 \label{eq:invariant}
\end{equation}
Thus a one-body energy measurement at fixed momentum cannot discriminate
these two proposed charge pairs at this order. It can test their common
conditional invariant. A charge or scattering-frame measurement can still
distinguish their interpretations. This algebra does not assume that either
pair has already been proved physical.

\begin{center}
\begin{tabular}{cccc}
\toprule
$\alpha$ & $p$ ($\epsilon=+1$) & $p_{\rm can}$ ($\epsilon=+1$) & $E_1|_P=-L_1$\\
\midrule
0.15 & $-0.157387066429$ & $-0.190035364892$ & $-0.001424902441$\\
0.50 & $-0.309798364829$ & $-0.603389892154$ & $-0.050091498992$\\
0.85 & $-0.311827472659$ & $-0.962723079255$ & $-0.223358039969$\\
\bottomrule
\end{tabular}
\end{center}
\Lc\ (\Ccheck) These are conditional coefficients of $\hbar$;
$E_1|_P$ is orientation independent. They are not additional determinant data.

\paragraph{Concrete Fermi-edge target, with qualifications. \Lc/\Ld.}
\Lc\ (\Ccheck) Analytic differentiation and a small-amplitude numerical
remainder check give
\begin{equation}
 L_1=\frac1{2\pi}\left[3\sin\alpha+
 \sum_{r=1}^2\sin\varphi_r(\cos\varphi_r-\cos\alpha)
                         \log|\tan(\varphi_r/2)|\right].
\end{equation}
The light-root logarithm cancels identically; this expression is analytic
at zero and has $L_1(0)=0$. Its derivatives give
\begin{align}
 p_{\rm can}&=-\frac{4\epsilon}{\pi}\alpha
       +\frac{8\epsilon}{9\pi}\alpha^3
       -\frac{8\epsilon}{45\pi}\alpha^5+O(\alpha^7),\\
 E_1|_P&=-\frac4{3\pi}\alpha^3+\frac{14}{45\pi}\alpha^5+O(\alpha^7),\\
 \lim_{\alpha\to0^+}\frac{E_1|_P}{|P_0|}&=-\frac2\pi
       =-0.63661977236758134308\ldots .\label{eq:target}
\end{align}
This is a sequential limit of the one-loop coefficient. A perturbative
comparison using the canonical momentum requires $\hbar/\alpha^2\ll1$;
using the mode label requires additionally the logarithm implicit in $p$ to
remain small relative to the classical momentum shift. \Ld\ Toda data could
test \eqref{eq:target} only after fixing the common physical momentum,
energy, $\hbar$ and good/bad continuation dictionary. This note does not
claim that dictionary, or a finite-$\hbar$ uniform edge limit. It does not
import any claimed Phase~2i coefficient.

\section{Two-soliton normalization: what is and is not fixed}
\paragraph{Continuum constant. \Lc\ (\Ccheck).}
For regular light solitons define the ordinary half-line integrals
\begin{equation}
 I_{12}=\frac1{4\pi}\int_{\R\setminus\{0\}}\theta_1\theta_2'\dd q,
 \qquad J_{12}=\tfrac12(I_{12}-I_{21}).
\end{equation}
The branch endpoints and integration by parts give
\begin{equation}
 \theta(0^\pm)=\pi(\pm1+\epsilon),\quad
 I_{12}+I_{21}=-\frac\pi2(\epsilon_1+\epsilon_2),\quad
 J_{12}-I_{12}=\frac\pi4(\epsilon_1+\epsilon_2).\label{eq:exchange}
\end{equation}
Independent quadratures at $(\alpha_1,\alpha_2)=(0.3,0.6)$ verify all four
orientations to better than $10^{-54}$. In particular $J_{++}=0.392706352268363$
and $J_{+-}=0.408542241291364$. These are continuum integrals, not an
absolutely normalized two-body quantum phase.

\paragraph{Conditional fixed-momentum expression. \Lc/\Ld.}
\Lc\ (\Ccheck, Taylor algebra) If a normalized determinant calculation
established $I_{12}$ as its label-frame one-loop phase up to a constant,
then changing from labels to physical momenta would give
\begin{equation}
 \delta_{1,P}=I_{12}+c_{12}
       -\sum_{j=1}^2p_j\frac{\partial\delta_0}{\partial P_{0j}},
 \qquad S=\exp\{i\delta_0/\hbar+i\delta_{1,P}+O(\hbar)\}.
 \label{eq:two}
\end{equation}
The $p_j$ must belong to that same determinant and state convention.
One cannot insert either charge pair by preference while retaining an
unconverted label-frame phase. \Ld\ Here the label-frame determinant
identification itself is not proved; thus the unknown is not yet known to
be only $c_{12}$. A boundary Schur contribution could change its derivatives.
\eqref{eq:two} isolates the constant only conditional on the FK-type gradient
being applicable to the completed boundary problem.

\paragraph{FK provenance and the role of $Z$.}
\La\ as recorded, unverified in this session: the repository's
\texttt{audit\_determinant\_foundations.tex} reports that FK \S5.4 explicitly
leaves a factor $Z$ to supplementary requirements and that its argument is
for a Klein--Gordon fluctuation problem. That audit records reading relevant
full-text passages, not the whole work. Neither FK full text nor its abstract
was read for this note; no new DOI, arXiv identifier or external theorem is
asserted. We use precisely the recorded limitation, not an unverified
extension of FK to Boussinesq.

\paragraph{Do unitarity, crossing and one-soliton limits fix it?}
\Lc\ (\Ccheck, elementary phase algebra) If label exchange has already been
identified with the inverse physical collision, scalar unitarity of
$Z_{12}e^{iI_{12}}$ requires
\begin{equation}
 Z_{12}Z_{21}=\exp\{i\pi(\epsilon_1+\epsilon_2)/2\}.
\end{equation}
Replacing $I$ by $J$ is one algebraic solution, leaving a factor
$e^{ic_{12}}$ with $c_{21}=-c_{12}$ modulo $2\pi$. It is not a determinant
calculation or a choice made here. In the same-direction sectors the right
side is $-1$; it is $+1$ head-on. Incoming/outgoing state conventions are
needed before applying this label test as physical unitarity.

\Lc\ (\Ccheck) As a precise limitation of bootstrap constraints alone,
consider a formal charge-reversal crossing map $(P,M)\mapsto(-P,-M)$.
The family
\begin{equation}
 F_{12}=\lambda(P_1M_2-P_2M_1),\qquad \lambda\in\R,
 \qquad S\mapsto e^{iF_{12}}S
\end{equation}
is antisymmetric, changes sign under crossing of either leg, and vanishes
when either leg has $P=M=0$. It therefore preserves unitarity, that inverse
crossing relation and the corresponding soft limits, while changing the
phase in the interior. This is an algebraic counterexample for the stated
map, not a derived Boussinesq crossing axiom. It would be excluded by a
fully established FK gradient; it is \emph{not} an extra constant consistent
with a fixed gradient.

\Ld\ There is no specified quantum Boussinesq crossing continuation or
normalized soft reference state in the supplied determinant construction.
If the gradient were established on a connected chamber and a normalized
one-soliton/soft limit were proved to approach a known phase there, that
reference would fix its remaining constant modulo $2\pi$. The limits and
continuation have not been established here, so these requirements presently
do not determine $Z$. Equal finite collective volumes, or two identical
asymptotic frequencies, do not replace the missing time-evolution phase.

\clearpage
\section{Check inventory, failures recorded, and the next required input}
\begin{longtable}{p{0.44\linewidth}p{0.12\linewidth}p{0.34\linewidth}}
\toprule Result & Label & Check/scope\\\midrule
Potential operator, boundary form, conserved charges, divergent tail norms
 & (c) & \G; exact algebra and 55-digit integrals\\
Bounded static kernel, velocity Jacobi field, continuum and heavy modes
 & (c) & \Scheck; exact identities, 60-digit residuals; no continuum completeness claim\\
Rank-two three-mode form, auxiliary completion, constraint and measure factors
 & (c) & \G; exact canonical embedding and independent field quadrature\\
Vacuum normal ordering and inherited regulated trace
 & (c) & \Ccheck; Wick algebra, 65-digit quadrature and cutoff asymptotics\\
Common fixed-$P$ coefficient and $-2/\pi$ target
 & (c), given trace & \Ccheck; derivative, Taylor and numerical remainder checks\\
Raw/antisymmetric phase constants and bootstrap limitation
 & (c) & \Ccheck; four orientations and symbolic charge algebra\\
Physical $L_1$, physical canonical charge, absolute $E(P)$ and $S$
 & (d) & No compatible full constrained determinant constructed\\
\bottomrule
\end{longtable}

\paragraph{Failures recorded.}
\begin{enumerate}
\item The three-tangent inverse-Wronskian construction fails because its
matrix is singular. This is a mathematical finding, not a numerical accident.
\item Treating shift and velocity/width as independent finite-norm
configuration collective paths fails by the linear box divergence
\eqref{eq:norm}. The auxiliary completion fixes the finite rank defect only.
\item The inherited heavy-mode shorthand assigns both growth rates to the
same spatial profile. The explicit residual quoted after \eqref{eq:heavy}
fails; the complementary profiles pass. No tolerance was weakened.
\item Dropping the equal-mode-count endpoint fails UV cancellation with the
specified vacuum counterterm. The endpoint was retained, not absorbed into
an adjustable subtraction.
\item During script development, SymPy structural equality of two equivalent
forms of the quartic second derivative failed. Testing their expanded
difference repaired the assertion; the formula and tolerance were unchanged.
\item During review of the script, a comment had the wrong sign in
$\pi_t=\mp\partial_x A$. It was corrected to $\pi_t=-\partial_x A$ and the
explicit symbolic conservation identity was added before the final run.
\item The first PDF build exposed an overfull heading and a missing row break
in the charge display. The heading, display and summary-table widths were
corrected; visual inspection also caught an orphaned table rule, repaired
by starting the check inventory on a new page before the final double build.
A subsequent text patch initially missed its context and was reapplied to
the actual paragraph; this did not change any calculation.
\end{enumerate}
The three scripts finish with \texttt{ALL CHECKS PASSED} for the above
checks only. JSON files with the matching \texttt{det\_...\_results.json}
names cache the residuals, conditional numbers, precision, scope and runtimes;
they are rewritten by a fresh run, not used to skip verification.
Latest measured runtimes on this host are 43.55 seconds (geometry),
1.91 seconds (spectrum), and 69.38 seconds (charges). Earlier runs were
approximately 22, 1.2 and 43 seconds respectively; wall times vary with
concurrent host load. The JSON files record the latest successful runs.

\paragraph{Precise remaining work and PI input. \Ld}
First specify which endpoint and shift-charge sectors define the quantum
one-soliton propagator, and whether comparisons across $M$ are boundary-state
comparisons or a constrained reduction with compensating fields. That is the
preference-dependent physical choice this note deliberately leaves to the PI.
Then construct a single finite spatial/temporal problem for field, constraints,
vacuum subtraction and normalized external states; retain its mixed blocks;
prove spectral completeness or compute its projected resolvent; and fix the
Gaussian square-root branch by evolution from a reference. This would test
whether the result is the conditional $L_1$, whether its velocity derivative
is the physical momentum, and only then whether the two-body remainder is
just FK's $Z$. The explicit heavy-branch instability and the cusp must also
be handled if those parameters are to represent external particles.
These are the exact missing steps, not a claimed no-go theorem.

\paragraph{Source record.}
Read locally: \texttt{AGENTS.md}, \texttt{README.md}, the ranked steps in
\texttt{report.tex}, and the relevant portions of
\texttt{oneloop\_phase.tex}, \texttt{oneloop.tex},
\texttt{audit\_determinant\_foundations.tex},
\texttt{audit\_projected\_measure.tex},
\texttt{audit\_charges\_action.tex}, and \texttt{charge\_frame\_master.tex},
together with the relevant phonon and charge scripts. No external paper
was read in this session. The only inherited literature statement used is
the explicitly qualified FK scope/constant statement above; its full-text
verification remains on the PI's source-access list. All new mathematical
claims are internal derivations, and no new bibliographic entry is invented.
\end{document}
